For a law \(P\) in the revised class, define the finite one-sided score measure
\[
 \mu_{t,P}(E)=P_X(E\cap\mathcal A_t).
\]
If \(x\in\mathcal B\) and \(r>0\), set \(s=\min\{h_0,r/2\}\). Positivity of \(m\) and \(\mathrm{ass:declared\text{-}modulus\text{-}lower\text{-}mass}\), which is a member property of \(\mathrm{def:modulus\text{-}profile\text{-}class}\), give
\[
 \mu_{t,P}\{z:d(z,x)<r\}\ge q_{t,P}(x,s)\ge m(s)>0. \tag{40}
\]
Thus \(x\in\mathcal X_{t,P}\) on both sides. The unique Lipschitz representative and the definition of its trace consequently give
\[
 \theta_{t,P}(x)=g_{t,P}(x). \tag{41}
\]
Under \(\mathcal A_0=\mathbb R^2\setminus\mathcal A_1\), sharp assignment makes \(X_i\in\mathcal A_t\) equivalent to \(T_i=t\). Therefore the conditional mean of each observed side average is the corresponding average of \(g_{t,P}(X_i)\). The bounded-trace member property yields \(\tau_P(x)\in[-2L,2L]\); projection onto this closed interval cannot increase pointwise distance to \(\tau_P(x)\), and the clipped loss is at most \(4L\).

We prove the upper display. Fix \(h\) with \(h,2h\in\mathcal H_n\) and \(nm(h)\ge A\log(e+n/h)\). Let \(z_1,\ldots,z_K\) be a maximal \(h\)-separated subset of \(\mathcal B\). It is an \(h\)-net, and \(\mathrm{ass:global\text{-}boundary\text{-}entropy}\) gives \(K\le C_{\mathcal B}h^{-1}\). For each \(j,t\), \(N_{t,n}(z_j,h)\) is binomial with success probability \(q_{t,P}(z_j,h)\ge m(h)\), by independent sampling and the class's lower-mass property. If \(N\sim\operatorname{Bin}(n,q)\), then Markov's inequality applied to \(2^{-N}\) gives
\[
 \Pr\{N\le nq/2\}
 \le 2^{nq/2}(1-q/2)^n
 \le \exp(-nq/8), \tag{42}
\]
where the last inequality follows from \(1-y\le e^{-y}\) and \((\log2)/2-1/2<-1/8\). Hence a union bound, monotonicity in the threshold, and \(q\ge m(h)\) show
\[
 \Pr_P(H_h^c)\le 2C_{\mathcal B}h^{-1}e^{-nm(h)/8}, \tag{43}
\]
where \(H_h\) is the event on which all \(N_{t,n}(z_j,h)\ge nm(h)/2\). For every \(x\in\mathcal B\), select a net point \(z_j\) with \(d(x,z_j)\le h\). The inclusion \(B(z_j,h)\subseteq B(x,2h)\) implies on \(H_h\) that
\[
 \inf_{x\in\mathcal B,\,t\in\{0,1\}}N_{t,n}(x,2h)\ge nm(h)/2. \tag{44}
\]
Choose the constant \(A\) so that the right-hand side of (44) is at least \(8\log(e+n/(2h))\). Then \(2h\in\widehat{\mathcal H}_{t,n}(x)\) for every \(x,t\) on \(H_h\).

Condition on the scores. For any selected nonempty averaging set at \((x,b,t)\), Lipschitz smoothness and (41) give
\[
 \left|\frac1{N_{t,n}(x,b)}\sum_{i:X_i\in\mathcal A_t,\ d(X_i,x)\le b}g_{t,P}(X_i)-\theta_{t,P}(x)\right|\le Lb. \tag{45}
\]
By \(\mathrm{lem:disk\text{-}pattern\text{-}gaussian\text{-}control}\), simultaneously over every candidate disk pattern, the conditional expected normalized Gaussian error is at most
\[
 \sqrt{2\log\!\left(4|\mathcal H_n|\sum_{j=0}^3\binom nj\right)}+1\le C\sqrt{\log(e+n)}. \tag{46}
\]
Because every \(b\in\mathcal H_n\) satisfies \(b\le h_0<1\), \(\log(e+n)\le C\log(e+n/b)\). Combining (45)--(46) with the minimizing property in \(\mathrm{def:count\text{-}adaptive\text{-}bandwidth}\), and comparing the selected criterion to the admissible candidate \(2h\), gives on \(H_h\)
\[
 \mathbb E_P\!\left[\sup_{x\in\mathcal B}|\widehat\tau_n(x)-\tau_P(x)|\mid X_1,\ldots,X_n\right]
 \le C\left\{Lh+\sigma\sqrt{\frac{\log(e+n/h)}{nm(h)}}\right\}. \tag{47}
\]
Here both side errors have been added, and (44) was used in the stochastic term. Projection leaves the error no larger. On \(H_h^c\), the clipped error is at most \(4L\). Equations (43) and (47) therefore yield
\[
 \sup_{P\in\mathcal P(m,\mathsf J_{-})}
 \mathbb E_P\sup_{x\in\mathcal B}|\overline\tau_n(x)-\tau_P(x)|
 \le C\left\{Lh+\sigma\sqrt{\frac{\log(e+n/h)}{nm(h)}}\right\}
   +CLh^{-1}e^{-nm(h)/8}. \tag{48}
\]
Increasing \(A\), if necessary, makes the last term at most \(CLh\), since \(e^{-nm(h)/8}\le(e+n/h)^{-A/8}\). The deterministic clipped bound supplies the first argument \(L\) of the minimum when no scale qualifies. Taking the infimum over qualifying dyadic pairs proves the upper bracket. The sample map in \(\mathrm{def:mass\text{-}adaptive\text{-}estimator}\) and the projection use neither \(m\) nor \(\mathsf J_{-}\).

We next prove the lower display within the same class. Choose any member \(P^\circ\) of the assumed nonempty class and retain its score marginal and assignment geometry. On an extension carrying independent standard normal variables \(Z_0,Z_1\), define a baseline law \(P_0\) by \(Y(t)=\sigma Z_t\). The score-dependent class properties are unchanged, while the side regressions and traces are zero, so the response smoothness, boundedness, Gaussian-error, sharp-assignment, and consistency properties all hold. Thus \(P_0\in\mathcal P(m,\mathsf J_{-})\).

Let \(M=\mathsf J_{-}(h,u)\ge1\). The class's lower-profile property and the definition of \(\mathsf J_P\) supply \(h\)-separated points \(x_1,\ldots,x_M\) for which \(a_{P_0}(x_j,h)\le u\). Choose \(t_j\) attaining this minimum and put
\[
 \eta_h=\frac{Lh}{12},\qquad
 \psi_{j,h}(z)=\left(1-\frac{3d(z,x_j)}h\right)_+.
\]
Define \(P_j\) by retaining the baseline score law and errors and shifting only the potential outcome on side \(t_j\): add \(\eta_h\psi_{j,h}(X)\) when \(t_j=1\), and subtract it when \(t_j=0\). The added regression is \(L/4\)-Lipschitz because \(\psi_{j,h}\) is \(3/h\)-Lipschitz; its absolute trace is at most \(Lh/12<L\), since \(h\le h_0<1\). All design properties remain unchanged, and the error around the shifted conditional mean is still \(\mathcal N(0,\sigma^2)\). Hence every \(P_j\) belongs to \(\mathcal P(m,\mathsf J_{-})\).

By (40)--(41), each center is in both one-sided supports and the target shift at \(x_j\) equals \(\eta_h\). If \(k\ne j\), then \(d(x_j,x_k)\ge h\), so \(\psi_{k,h}(x_j)=0\). Consequently
\[
 \|\tau_{P_j}-\tau_{P_k}\|_\infty\ge\eta_h,
 \qquad 0\le j<k\le M. \tag{49}
\]
Conditionally on a score value on the perturbed assignment side, the observed outcome changes from a normal mean zero to a normal mean \(s_j\eta_h\psi_{j,h}(X)\), with \(s_j\in\{-1,1\}\), and it does not change elsewhere. Direct integration of the Gaussian log-likelihood ratio gives conditional relative entropy \(\eta_h^2\psi_{j,h}(X)^2/(2\sigma^2)\). Since the tent is supported inside \(B(x_j,h/3)\subseteq B(x_j,h)\),
\[
 \operatorname{KL}(P_j^{\otimes n},P_0^{\otimes n})
 =\frac{n\eta_h^2}{2\sigma^2}\int\psi_{j,h}(z)^2\,d\mu_{t_j,P_0}(z)
 \le \frac{nL^2h^2u}{288\sigma^2}. \tag{50}
\]
If the first information condition holds, retain any one alternative and apply the two-law clause of \(\mathrm{lem:finite\text{-}information\text{-}testing}\), decreasing the numerical \(a\) so that (50) is at most \(1/8\). If \(M\ge2\) and the logarithmic condition holds, choose \(a\) so small that (50) is at most the multiple-law threshold times \(\log M\); this is possible uniformly over \(M\ge2\) because \(\log(M+1)/\log M\le\log3/\log2\). Applying the same lemma with the sup metric and \(2\delta=\eta_h\) proves \(R_n\{\mathcal P(m,\mathsf J_{-})\}\ge cLh\).

For a sequence of nonempty classes, suppose \(h_n\downarrow0\), \(h_n,2h_n\in\mathcal H_n\), and
\[
 \frac{nm_n(h_n)}{\log(e+n/h_n)}\longrightarrow\infty. \tag{51}
\]
Then the eligibility condition holds eventually, \(Lh_n\to0\), and the square-root term in (48) tends to zero. This proves sufficient uniform consistency. Conversely, if the supremum of the lower-bound-feasible radii does not tend to zero, then along a subsequence one may choose feasible radii bounded below by a fixed positive constant (or half that constant when the supremum is not attained). The lower display then keeps the minimax risk bounded below by a positive constant, proving the stated necessary condition.

For \(m_\alpha\), \(\mathrm{prop:exponential\text{-}horn\text{-}consistency}\) supplies positivity and
\(m_\alpha(h)\asymp h^{\alpha+1}e^{-h^{-\alpha}}\). Set \(s_n=\{(\log n)/2\}^{-1/\alpha}\). Then
\[
 \frac{nm_\alpha(s_n)}{\log(e+n/s_n)}
 \asymp
 \frac{n^{1/2}s_n^{\alpha+1}}{\log n}\longrightarrow\infty. \tag{52}
\]
Round \(s_n\) upward to the next point \(h_n\) of the dyadic grid. Eventually \(s_n\le h_n<2s_n\), both \(h_n\) and \(2h_n\) belong to \(\mathcal H_n\), and monotonicity of \(m_\alpha\) preserves (52). Thus the sufficient condition holds for the fixed exponential horn.

For polynomial profiles, set \(u=C_mh^\kappa\). If \(M\asymp1\), (50) is bounded at \(h\asymp n^{-1/(\kappa+2)}\). If \(M\asymp h^{-1}\), the logarithmic information condition is met at \(h\asymp(\log n/n)^{1/(\kappa+2)}\), because \(nh^{\kappa+2}\asymp\log(1/h)\). The matching achievability bounds in \(\mathrm{thm:thickness\text{-}entropy\text{-}separation}\) make both rates exact.

Finally, the two generic displays can differ because the declared lower packing profile may be slack while the upper proof covers the entire boundary. On the circular pervasive witness, take \(m(h)=c_mh^\kappa\) and declare \(\mathsf J_{-}(h,u)=0\) for \(u<C_mh^\kappa\) and \(\mathsf J_{-}(h,u)=1\) for \(u\ge C_mh^\kappa\). By \(\mathrm{lem:circular\text{-}witness\text{-}membership}\), after choosing its fixed constants this defines a nonempty class containing the pervasive witness. Its declared lower profile exposes only \(M=1\), whereas (48) uses the global boundary entropy and therefore carries the global logarithm. This legal declaration certifies the stated possible gap and completes the two-sided bracket.